% ============================================================
%  main.tex  —  AI Tools for Research
%  Metropolis Beamer Presentation
%  Compile with: latexmk  (see latexmkrc)
% ============================================================

\documentclass[aspectratio=169, 10pt]{beamer}

% ---------- Theme -------------------------------------------
\usetheme[progressbar=frametitle, block=fill, numbering=fraction]{metropolis}

% ---------- Fonts -------------------------------------------
\usepackage[T1]{fontenc}
\usepackage[utf8]{inputenc}
\usepackage{FiraSans}
\usepackage{FiraMono}

% ---------- Math --------------------------------------------
\usepackage{amsmath, amssymb, mathtools}

% ---------- Code listings -----------------------------------
\usepackage{listings}
\lstset{
  basicstyle=\ttfamily\footnotesize,
  breaklines=true,
  breakatwhitespace=true,
  keywordstyle=\color{mLightBrown}\bfseries,
  commentstyle=\color{mDarkTeal}\itshape,
  stringstyle=\color{mLightGreen},
  numbers=left,
  numberstyle=\tiny\color{gray},
  numbersep=5pt,
  frame=single,
  framerule=0.4pt,
  framesep=3pt,
  xleftmargin=12pt,
  xrightmargin=4pt,
  showstringspaces=false,
  tabsize=2,
}
\lstdefinelanguage{json}{
  morestring=[b]",
  morecomment=[l]{//},
  morecomment=[s]{/*}{*/},
  morekeywords={true,false,null},
  sensitive=false,
}
\lstdefinelanguage{python}{
  keywords={def,class,import,from,return,if,else,elif,for,while,
            in,not,and,or,True,False,None,with,as,try,except,raise,
            lambda,yield,pass,break,continue},
  morestring=[b]",
  morestring=[b]',
  morecomment=[l]{\#},
  sensitive=true,
}
\lstdefinelanguage{bash}{
  keywords={ls,cd,mkdir,cp,mv,rm,echo,cat,grep,find,make,git,
            pip,python,conda,curl,wget},
  morecomment=[l]{\#},
  morestring=[b]",
  morestring=[b]',
  sensitive=false,
}

% ---------- TikZ / pgfplots --------------------------------
\usepackage{tikz}
\usetikzlibrary{arrows.meta, positioning, shapes.geometric, fit,
                backgrounds, calc, decorations.pathreplacing}
\usepackage{pgfplots}
\pgfplotsset{compat=1.18}

% ---------- Tables -----------------------------------------
\usepackage{booktabs, tabularx, multirow, colortbl}
\newcolumntype{Y}{>{\centering\arraybackslash}X}

% ---------- Graphics ---------------------------------------
\usepackage{graphicx}

% ---------- Bibliography -----------------------------------
\usepackage[backend=biber, style=authoryear-comp, maxcitenames=2,
            uniquelist=false, uniquename=false]{biblatex}
\addbibresource{refs.bib}

% ---------- Utilities --------------------------------------
\usepackage{appendixnumberbeamer}
\usepackage{qrcode}
\usepackage{hyperref}
\usepackage[normalem]{ulem}
\hypersetup{colorlinks=true, linkcolor=., urlcolor=mLightBrown}

% ---------- Custom commands --------------------------------
% Note: mLightBrown, mDarkTeal, mLightGreen are defined by metropolis theme
\newcommand{\tool}[1]{\textbf{\textcolor{mLightBrown}{#1}}}
\newcommand{\model}[1]{\texttt{#1}}

% Demo placeholder: standout frame with large "DEMO" label
\newcommand{\demoframe}[1]{%
  \begin{frame}[standout]
    \Huge $\blacktriangleright$\\[0.5em]
    {\large\textsc{Live Demo}}\\[1em]
    \normalsize #1
  \end{frame}%
}

% ---------- Metadata ---------------------------------------
\title{AI Tools for Research}
\subtitle{A Practical Guide for CS/AI Researchers}
\author{%
  David Adams\\
  \small\href{mailto:david.adams1@student.unimelb.edu.au}{david.adams1@student.unimelb.edu.au}%
}
\institute{FEIT \\ The University of Melbourne}
\date{23 June 2026}

% ============================================================
\begin{document}
% ============================================================

\maketitle

% ---- About This Talk --------------------------------------
\begin{frame}[shrink=5]{About This Talk}
  \begin{columns}[T, onlytextwidth]
    \begin{column}{0.48\textwidth}
      \textbf{What we cover}
      \begin{itemize}\setlength{\itemsep}{1pt}
        \item AI in software development broadly
        \item The frontier model landscape
        \item Chat \& reasoning workflows
        \item Coding assistants for researchers
        \item Agentic \& MCP pipelines
        \item Literature discovery tools
        \item Writing assistance
        \item Caveats, ethics, best practices
      \end{itemize}
    \end{column}
    \hfill
    \begin{column}{0.48\textwidth}
      \textbf{What we do \emph{not} cover}
      \begin{itemize}\setlength{\itemsep}{1pt}
        \item Model internals / training details
        \item Fine-tuning or RLHF
        \item Deployment / MLOps
        \item Robotics / embodied AI
        \item A comprehensive product review (landscape moves fast)
      \end{itemize}
      \vspace{0.4em}
      \textbf{Format:} Talk + live demos\\
      \textbf{Slides:} Posted after the talk
    \end{column}
  \end{columns}
\end{frame}

% ---- Table of Contents ------------------------------------
\begin{frame}{Outline}
  \tableofcontents
\end{frame}

% ===========================================================
%  Sections
% ===========================================================
% ============================================================
%  Section 1: AI in Software Development (Broad)
% ============================================================

\section{AI in Software Development}

% ---- Timeline slide ----------------------------------------
\begin{frame}{The Shift in How Code Gets Written}
  \centering
  \begin{tikzpicture}[
    event/.style={draw, rounded corners=2pt, fill=mLightBrown!15,
                  minimum width=2.6cm, minimum height=1.0cm,
                  align=center, font=\small},
    arr/.style={-{Stealth[length=5pt]}, thick, mLightBrown},
  ]
    \node[event] (e1) at (0,0)    {2021\\IntelliCode\\autocomplete};
    \node[event] (e2) at (3.4,0)  {2022\\GitHub Copilot\\GA release};
    \node[event] (e3) at (6.8,0)  {2023\\ChatGPT boom\\chat-with-code};
    \node[event] (e4) at (10.2,0) {2024--26\\LLM-native IDEs\\agentic coding};

    \draw[arr] (e1.east) -- (e2.west);
    \draw[arr] (e2.east) -- (e3.west);
    \draw[arr] (e3.east) -- (e4.west);

    % Timeline bar
    \draw[thick, gray!50] (-1.5,-0.7) -- (12.0,-0.7);
  \end{tikzpicture}

  \vspace{2.5em}
  \begin{itemize}
    \item From \textbf{single-line autocomplete} → whole-function generation → full-feature agents
    \item Paradigm shift: the developer describes \emph{intent}, the model produces \emph{code}
    \item Researcher impact: prototyping, baselines, data pipelines all affected
  \end{itemize}
\end{frame}

% ---- Ecosystem slide ----------------------------------------
\begin{frame}[shrink=5]{The Developer AI Tool Ecosystem}
  \begin{columns}[T, onlytextwidth]
    \begin{column}{0.32\textwidth}
      \textbf{Standalone Applications}\\[0.3em]
      \begin{itemize}
        \item \tool{GitHub Copilot} (VS Code, JetBrains, Vim)
        \item \tool{Cursor} (fork of VS Code)
        \item \tool{Codex} (agentic interface)
        \item \tool{Claude Cowork} (chat and computer use)
        \item \tool{Google Antigravity (v1/v2)} (agentic dev platform)
        \item \sout{\tool{Continue}} (aquired by cursor)
      \end{itemize}
    \end{column}
    \hfill
    \begin{column}{0.32\textwidth}
      \textbf{Terminal / Agentic}\\[0.3em]
      \begin{itemize}
        \item \tool{Claude Code} (Anthropic CLI)
        \item \tool{Codex} (OpenAI CLI)
        \item \tool{agy} (Google CLI)
        \item \tool{OpenCode} (open-source terminal/IDE/desktop agent; 75+ LLM providers)
        \item \tool{Pi} (minimal terminal agent)
      \end{itemize}
    \end{column}
    \hfill
    \begin{column}{0.32\textwidth}
      \textbf{Code Review / PR}\\[0.3em]
      \begin{itemize}
        \item \tool{GitHub Copilot} PR review
        \item \tool{Sourcery} (refactoring)
        \item \tool{CodeRabbit} (AI PR reviews)
        \item \tool{Sweep} (issue-to-PR agent)
      \end{itemize}
    \end{column}
  \end{columns}
\end{frame}

% ---- What these tools do -----------------------------------
\begin{frame}{What These Tools Actually Do}
  \begin{columns}[T, onlytextwidth]
    \begin{column}{0.30\textwidth}
      \begin{block}{Completion}
        \begin{itemize}
          \item Next-line / next-block autocomplete
          \item Triggered by context in the cursor
          \item Very low latency
          \item \textit{Examples:} Copilot inline, Tabnine
        \end{itemize}
      \end{block}
    \end{column}
    \hfill
    \begin{column}{0.30\textwidth}
      \begin{block}{Chat-with-Codebase}
        \begin{itemize}
          \item Conversational interface with codebase context
          \item Ask questions, get explanations
          \item Edit specific files on request
          \item \textit{Examples:} Cursor, OpenCode, Continue
        \end{itemize}
      \end{block}
    \end{column}
    \hfill
    \begin{column}{0.30\textwidth}
      \begin{block}{Agentic}
        \begin{itemize}
          \item Multi-step task execution
          \item File R/W, shell commands, git
          \item Iterates until task complete
          \item \textit{Examples:} Claude Code, Devin
        \end{itemize}
      \end{block}
    \end{column}
  \end{columns}
\end{frame}

% ---- Adoption signals --------------------------------------
\begin{frame}[shrink=5]{Adoption Is High; Trust Is the Bottleneck}
  \begin{columns}[T, onlytextwidth]
    \begin{column}{0.48\textwidth}
      \textbf{Stack Overflow Survey 2025}
      \begin{itemize}\setlength{\itemsep}{2pt}
        \item \textbf{84\%} of all respondents are using or planning to use AI tools, rising from 76\% in previous surveys.
        \item \textbf{51\% of professional developers} now interact with AI coding assistants daily.
        \item \alert{Trust is falling}: Only \textbf{29\% of developers} actually trust AI outputs in 2025, down 11 percentage points from 2024.
      \end{itemize}
    \end{column}
    \hfill
    \begin{column}{0.48\textwidth}
      \textbf{GitHub Octoverse 2025}
      \begin{itemize}\setlength{\itemsep}{2pt}
        \item Over \textbf{1.1 million public repositories} on GitHub used or integrated with an LLM SDK.
        \item \textbf{80\% of new developers} choose to use GitHub Copilot during their first week of onboarding.
      \end{itemize}
      \vspace{0.4em}
      \begin{block}{TLDR;}
        The initial developer hype cycle has settled into a pragmatism gap: massive daily deployment is mismatched with low trust in silent errors.
      \end{block}
    \end{column}
  \end{columns}
  \vspace{0.6em}
  \centering\tiny\textit{Sources: Stack Overflow Developer Survey 2025, GitHub Octoverse 2025 Reports.}
\end{frame}

% ---- AI coding products table ------------------------------
\begin{frame}{AI Coding Tools: Product Landscape I}
\centering
\scriptsize
\begin{tabular}{p{2.6cm} p{2.7cm} c c p{1.7cm} p{1.4cm}}
\toprule
\textbf{Product} & \textbf{Type} & \textbf{Open} & \textbf{BYOM} & \textbf{Price} & \textbf{Released} \\
\midrule
Cursor & Standalone IDE; Local/CLI; Cloud & -- & -- & \$0--200/mo & Mar 2023 \\
Claude Code & IDE Ext.; Local/CLI; Cloud & -- & -- & \$20--200/mo & Feb 2025 \\
GitHub Copilot Coding Agent & IDE Ext.; Cloud & -- & -- & \$0--39/mo & May 2025 \\
Windsurf & Standalone IDE & -- & -- & \$0--200/mo & Nov 2024 \\
Codex & IDE Ext.; Local/CLI; Cloud & -- & -- & \$0--200/mo & Apr 2025 \\
Cline & IDE Ext.; Local/CLI & -- & -- & Free + API & Jul 2024 \\
Devin & Cloud & -- & -- & \$20--500/mo & Mar 2024 \\
Aider & Local/CLI & -- & -- & Free + API & Jun 2023 \\
Google Antigravity & Standalone IDE & -- & -- & \$0--250/mo & Nov 2025 \\
Roo Code & IDE Extension & -- & -- & Free + API & Nov 2024 \\
Amazon Q Developer & IDE Ext.; Local/CLI & -- & -- & \$0--19/mo & Apr 2024 \\
Gemini Code Assist & IDE Ext.; Local/CLI & -- & -- & \$0--45+/mo & Apr 2024 \\
Continue & IDE Ext.; Local/CLI & Yes & Yes & \$0--20/mo & Jul 2023 \\
OpenHands & Cloud & Yes & Yes & \$0--500/mo & Mar 2024 \\
Zed & Standalone IDE & -- & -- & \$0--10/mo & Jan 2024 \\
\bottomrule
\end{tabular}
\end{frame}

\begin{frame}{AI Coding Tools: Product Landscape II}
\centering
\scriptsize
\begin{tabular}{p{2.6cm} p{2.7cm} c c p{1.7cm} p{1.4cm}}
\toprule
\textbf{Product} & \textbf{Type} & \textbf{Open} & \textbf{BYOM} & \textbf{Price} & \textbf{Released} \\
\midrule
Augment Code & IDE Ext.; Local/CLI & -- & -- & \$20--200/mo & Apr 2024 \\
Amp & IDE Ext.; Local/CLI & -- & -- & Free + API & May 2025 \\
Kiro & Standalone IDE; Local/CLI & -- & -- & \$0--200/mo & Jul 2025 \\
Jules & Cloud & -- & -- & \$0--250/mo & May 2025 \\
Gemini CLI & Local/CLI & -- & Yes & Free + API & Jun 2025 \\
Goose & Local/CLI & Yes & Yes & Free + API & Jan 2025 \\
Warp 2.0 & Local/CLI & -- & -- & \$0--200/mo & Jun 2025 \\
Genie & Cloud & -- & -- & \$0--99/mo & Aug 2024 \\
Manus & Cloud & -- & -- & \$0--199/mo & Mar 2025 \\
Mistral Vibe & IDE Ext.; Local/CLI & -- & -- & \$0--50/mo & Jun 2025 \\
opencode & IDE Ext.; Local/CLI & Yes & Yes & Free + API & Jun 2025 \\
Qoder & Standalone IDE & -- & -- & \$0--100/mo & Aug 2025 \\
Qwen Code & Local/CLI & Yes & Yes & Free + API & Jul 2025 \\
Kimi CLI & IDE Ext.; Local/CLI & -- & Yes & Free + API & Oct 2025 \\
BLACKBOX & IDE Ext.; IDE; CLI; Cloud & -- & -- & \$10--40/mo & Sep 2022 \\
\bottomrule
\end{tabular}
\end{frame}

% ---- Demo placeholder -------------------------------------
\demoframe{GitHub Copilot / Cursor — in-editor tab completion \& chat on a real snippet}

% ============================================================
%  Section 2: The Model Landscape
% ============================================================

\section{The Model Landscape}

% ---- Frontier players table --------------------------------
\begin{frame}{Frontier Players}
  \small
  \begin{tabularx}{\textwidth}{lYYYY}
    \toprule
    \textbf{Model} & \textbf{Context window} & \textbf{Multimodal} & \textbf{API access} & \textbf{Open weight} \\
    \midrule
    \tool{Claude 3.5 / Opus 4} & 200K & Vision & Yes & No \\
    \tool{GPT-4o}              & 128K & Vision + Audio & Yes & No \\
    \tool{Gemini 1.5 Pro}      & 1M   & Vision + Audio & Yes & No \\
    \tool{Llama 3.1 405B}      & 128K & Text & Via providers & \checkmark \\
    \tool{Mistral Large}       & 128K & Text & Yes & Partial \\
    \tool{Qwen2.5-72B}         & 128K & Vision & Via providers & \checkmark \\
    \tool{DeepSeek-R1}         & 64K  & Text & Yes & \checkmark \\
    \bottomrule
  \end{tabularx}
  \vspace{0.3em}

  \begin{itemize}
    \item \textbf{Capability gap is shrinking} — open-weight models are increasingly competitive
    \item Context window size matters for long papers, large codebases, and multi-document tasks
  \end{itemize}
\end{frame}

% ---- Capability tiers --------------------------------------
\begin{frame}{Capability Tiers}
  \begin{columns}[T, onlytextwidth]
    \begin{column}{0.48\textwidth}
      \begin{block}{Chat / Instruction Models}
        \begin{itemize}
          \item Fast, conversational, general-purpose
          \item Strong instruction following
          \item \textbf{Best for:} writing, summarization,\newline Q\&A, coding assistance
          \item \textit{Examples:} \model{claude-3.5-sonnet},\newline \model{gpt-4o}, \model{gemini-flash}
        \end{itemize}
      \end{block}
    \end{column}
    \hfill
    \begin{column}{0.48\textwidth}
      \begin{block}{Reasoning / Thinking Models}
        \begin{itemize}
          \item Slower; internally ``thinks'' before responding
          \item Better at multi-step logical problems
          \item \textbf{Best for:} math proofs, algorithm design,\newline complex debugging, formal reasoning
          \item \textit{Examples:} \model{o1}, \model{o3}, \model{claude-3.7-sonnet} (extended thinking), \model{deepseek-r1}
        \end{itemize}
      \end{block}
    \end{column}
  \end{columns}
  \vspace{0.6em}
  \alert{Rule of thumb:} Start with a fast model; escalate to a reasoning model when the problem requires structured multi-step inference.
\end{frame}

% ---- Benchmarks slide --------------------------------------
\begin{frame}[shrink=12]{Benchmarks — Useful Signal, Not Ground Truth}
  \begin{columns}[T, onlytextwidth]
    \begin{column}{0.55\textwidth}
      \begin{tikzpicture}
        \begin{axis}[
          width=\textwidth, height=4.2cm,
          xbar, bar width=0.35cm,
          symbolic y coords={SWEBench, HumanEval, MMLU},
          ytick=data,
          xmin=0, xmax=105,
          xlabel={Score (\%)},
          nodes near coords,
          nodes near coords align={horizontal},
          every node near coord/.append style={font=\tiny},
          axis lines*=left,
          xmajorgrids=true,
          grid style={dashed, gray!30},
          legend style={at={(0.98,0.05)}, anchor=south east, font=\tiny},
          label style={font=\small},
          tick label style={font=\small},
        ]
          \addplot[fill=mLightBrown!70] coordinates {
            (49,SWEBench) (92,HumanEval) (88,MMLU)
          };
          \addplot[fill=blue!50] coordinates {
            (49,SWEBench) (90,HumanEval) (87,MMLU)
          };
          \addplot[fill=teal!50] coordinates {
            (35,SWEBench) (74,HumanEval) (80,MMLU)
          };
          \legend{Claude family, GPT-4o, Llama 3.1}
        \end{axis}
      \end{tikzpicture}
      \tiny Approximate/illustrative figures; see official leaderboards.
    \end{column}
    \hfill
    \begin{column}{0.42\textwidth}
      \textbf{Key benchmarks}\\[0.3em]
      \begin{itemize}
        \item \textbf{MMLU}~\parencite{hendrycks2021mmlu}: broad knowledge (57 subjects)
        \item \textbf{HumanEval}~\parencite{chen2021codex}: Python function synthesis
        \item \textbf{SWE-Bench}~\parencite{jimenez2024swebench}: real GitHub issue resolution
      \end{itemize}
      \vspace{0.6em}
      \textbf{Caveats}\\[0.3em]
      \begin{itemize}
        \item Contamination risk — benchmarks leak into training data~\parencite{srivastava2022bigbench}
        \item Does not capture latency, cost, or UX
        \item Task-specific evals often more informative
      \end{itemize}
    \end{column}
  \end{columns}
\end{frame}

% ---- Open-weight models ------------------------------------
\begin{frame}[fragile]{Open-Weight Models}
  \begin{columns}[T, onlytextwidth]
    \begin{column}{0.55\textwidth}
      \textbf{Leading families (2024--25)}\\[0.4em]
      \begin{itemize}
        \item \tool{Llama 3.1}~\parencite{touvron2023llama2}: 8B / 70B / 405B — Meta's workhorse
        \item \tool{Mistral}~\parencite{jiang2023mistral}: efficient, strong at code; Mixtral MoE variant
        \item \tool{Qwen 2.5}: Alibaba; excellent multilingual \& coding
        \item \tool{DeepSeek-R1}: reasoning model released with full weights
        \item \tool{Gemma 2}: Google's research-friendly release
      \end{itemize}
    \end{column}
    \hfill
    \begin{column}{0.42\textwidth}
      \textbf{Local deployment}\\[0.3em]
      Run any open-weight model locally with \tool{Ollama}:
      \begin{lstlisting}[language=bash, numbers=none]
# install and pull
ollama pull llama3.1:70b

# run inference
ollama run llama3.1:70b
      \end{lstlisting}
      \vspace{0.4em}
      \begin{itemize}
        \item \textbf{Use case:} sensitive data, airgapped environments, no API cost
        \item Also: \tool{vLLM}, \tool{llama.cpp}, \tool{LM Studio}
      \end{itemize}
    \end{column}
  \end{columns}
\end{frame}

% ---- Choosing a model — decision tree ----------------------
\begin{frame}[shrink=15]{Choosing a Model for Your Task}
  \centering
  \scalebox{0.85}{%
  \begin{tikzpicture}[
    decision/.style={diamond, draw, fill=mLightBrown!20, aspect=2.5,
                     align=center, font=\tiny, inner sep=1pt},
    result/.style={draw, rounded corners=2pt, fill=blue!10,
                   align=center, font=\tiny, minimum width=2.0cm,
                   minimum height=0.6cm},
    arr/.style={-{Stealth[length=4pt]}, thick},
    lbl/.style={font=\tiny, midway},
  ]
    \node[decision] (q1) {Sensitive / private\\data?};
    \node[result, left=1.6cm of q1]  (local) {Local open-weight\\\tool{Ollama}/\tool{vLLM}};
    \node[decision, below=0.75cm of q1] (q2) {Needs long context\\(\textgreater 64K tokens)?};
    \node[result, right=1.6cm of q2] (gemini) {\tool{Gemini 1.5 Pro}\\\tool{Claude} (200K)};
    \node[decision, below=0.75cm of q2] (q3) {Multi-step reasoning\\ / formal proof?};
    \node[result, right=1.6cm of q3] (reason) {\tool{o1}/\tool{o3}\\\tool{Claude + Ext. Think.}};
    \node[decision, below=0.75cm of q3] (q4) {Cost-sensitive\\ or high volume?};
    \node[result, right=1.6cm of q4] (cheap) {\tool{Gemini Flash}\\\tool{Claude Haiku}};
    \node[result, below=0.7cm of q4]  (general) {\tool{Claude Sonnet}\\\tool{GPT-4o} (general use)};

    \draw[arr] (q1.west)  -- (local.east)   node[lbl, above]{Yes};
    \draw[arr] (q1.south) -- (q2.north)     node[lbl, right]{No};
    \draw[arr] (q2.east)  -- (gemini.west)  node[lbl, above]{Yes};
    \draw[arr] (q2.south) -- (q3.north)     node[lbl, right]{No};
    \draw[arr] (q3.east)  -- (reason.west)  node[lbl, above]{Yes};
    \draw[arr] (q3.south) -- (q4.north)     node[lbl, right]{No};
    \draw[arr] (q4.east)  -- (cheap.west)   node[lbl, above]{Yes};
    \draw[arr] (q4.south) -- (general.north)node[lbl, right]{No};
  \end{tikzpicture}}
\end{frame}

% ---- Demo placeholder -----------------------------------------------
\demoframe{Side-by-side model comparison — same research prompt in Claude, GPT-4o, and Llama}

% ============================================================
%  Section 3: Chat and Reasoning Workflows for Research
% ============================================================

\section{Chat \& Reasoning Workflows}

% ---- Research prompt loop ----------------------------------
\begin{frame}{The Research Prompt Loop}
  \centering
  \begin{tikzpicture}[
    node distance=1.5cm and 2.2cm,
    box/.style={draw, rounded corners=4pt, minimum width=2.2cm,
                minimum height=0.85cm, align=center, font=\small,
                fill=mLightBrown!15},
    arr/.style={-{Stealth[length=5pt]}, thick, mLightBrown},
  ]
    \node[box] (res)    {Researcher};
    \node[box, right=of res] (prompt) {Prompt};
    \node[box, right=of prompt] (model) {Model};
    \node[box, below=of model] (output) {Response};
    \node[box, left=of output] (refine) {Refine\\/ Critique};
    \node[box, left=of refine] (insight) {Insight\\\& Action};

    \draw[arr] (res.east)    -- (prompt.west);
    \draw[arr] (prompt.east) -- (model.west);
    \draw[arr] (model.south) -- (output.north);
    \draw[arr] (output.west) -- (refine.east);
    \draw[arr] (refine.west) -- (insight.east);
    \draw[arr] (insight.north) |- (res.south);
  \end{tikzpicture}

  \vspace{0.7em}
  \begin{itemize}
    \item The loop is \textbf{iterative}: outputs inform the next prompt
    \item Quality of prompting directly determines quality of outputs
    \item Build a personal library of effective prompt patterns
  \end{itemize}
\end{frame}

% ---- Daily research use cases ------------------------------
\begin{frame}{Use Cases in Daily Research}
  \begin{columns}[T, onlytextwidth]
    \begin{column}{0.48\textwidth}
      \textbf{Paper comprehension}\\
      \begin{itemize}
        \item Paste abstract + intro, ask for ELI-researcher
        \item ``What are the key assumptions?''
        \item ``Compare to [prior work X]''
        \item Identify methodological gaps
      \end{itemize}
      \vspace{0.6em}
      \textbf{Counter-argument generation}\\
      \begin{itemize}
        \item ``Act as a skeptical reviewer. What are the three strongest objections to this claim?''
        \item Stress-test your own hypotheses before submission
      \end{itemize}
    \end{column}
    \hfill
    \begin{column}{0.48\textwidth}
      \textbf{Brainstorming}\\
      \begin{itemize}
        \item Hypothesis generation at low cost
        \item Experiment design brainstorming
        \item Related-work scoping
      \end{itemize}
      \vspace{0.6em}
      \textbf{Sanity-checking}\\
      \begin{itemize}
        \item ``Does this derivation look correct?''
        \item ``What edge cases am I missing?''
        \item Quick literature cross-check (verify independently!)
      \end{itemize}
    \end{column}
  \end{columns}
\end{frame}

% ---- Prompting patterns ------------------------------------
\begin{frame}{Prompting Patterns That Work}
  \begin{columns}[T, onlytextwidth]
    \begin{column}{0.48\textwidth}
      \begin{block}{Role Prompting}
        \small\textit{``You are an expert in [domain] with 20 years of experience. Review the following methodology and identify weaknesses.''}\\[0.2em]
        Sets tone, terminology, and depth of response.
      \end{block}
      \vspace{0.4em}
      \begin{block}{Decomposition}
        \small\textit{``Break this problem into 3--5 independent subproblems. Solve each, then synthesize.''}\\[0.2em]
        Reduces hallucination on complex tasks.~\parencite{wei2022cot}
      \end{block}
    \end{column}
    \hfill
    \begin{column}{0.48\textwidth}
      \begin{block}{Chain-of-Thought}
        \small\textit{``Think step by step. Show your reasoning before giving the final answer.''}\\[0.2em]
        Dramatically improves multi-step math and logic~\parencite{wei2022cot}.
      \end{block}
      \vspace{0.4em}
      \begin{block}{Critique \& Revise}
        \small\textit{``Here is a draft. Identify any factual errors, logical gaps, or unclear sections. Then produce a revised version.''}\\[0.2em]
        Two-pass approach improves quality.
      \end{block}
    \end{column}
  \end{columns}
\end{frame}

% ---- Reasoning models --------------------------------------
\begin{frame}{Reasoning Models for Hard Problems}
  \begin{columns}[T, onlytextwidth]
    \begin{column}{0.55\textwidth}
      \textbf{What makes them different}\\[0.3em]
      \begin{itemize}
        \item Internal ``thinking'' chain before the final answer
        \item Optimized for multi-step inference, not raw speed
        \item \model{o1} / \model{o3}: OpenAI's chain-of-thought-at-inference models
        \item \model{claude-3.7-sonnet} extended thinking: configurable thinking budget
        \item \model{deepseek-r1}: open-weight reasoning model with visible CoT
      \end{itemize}
    \end{column}
    \hfill
    \begin{column}{0.42\textwidth}
      \textbf{When to use them}\\[0.3em]
      \begin{itemize}
        \item Mathematical proofs and derivations
        \item Algorithm design and complexity analysis
        \item Debugging subtle logical errors
        \item Multi-constraint optimization problems
        \item Long-range planning tasks
      \end{itemize}
      \vspace{0.5em}
      \alert{Cost:} Reasoning models are \textbf{3--10$\times$ more expensive} and slower. Use them selectively.
    \end{column}
  \end{columns}
\end{frame}

% ---- Context window as a research tool ---------------------
\begin{frame}[shrink=8]{Context Window as a Research Tool}
  \begin{columns}[T, onlytextwidth]
    \begin{column}{0.55\textwidth}
      \textbf{What large contexts enable}\\[0.3em]
      \begin{itemize}
        \item Paste entire papers (PDFs $\approx$ 8K--40K tokens)
        \item Load full codebases into context
        \item Multi-document synthesis and comparison
        \item Long conversation history without truncation
      \end{itemize}
      \vspace{0.5em}
      \textbf{Current limits}\\[0.3em]
      \begin{itemize}
        \item GPT-4o: 128K tokens $\approx$ 4--5 long papers
        \item Claude: 200K tokens $\approx$ 7--8 long papers
        \item Gemini 1.5 Pro: 1M tokens $\approx$ a whole library section
      \end{itemize}
    \end{column}
    \hfill
    \begin{column}{0.42\textwidth}
      \begin{alertblock}{``Lost in the Middle''}
        Models perform worse on information embedded in the \emph{middle} of very long contexts~\parencite{liu2023lostmiddle}.\\[0.5em]
        \textbf{Mitigations:}
        \begin{itemize}
          \item Place critical information at \textbf{start or end}
          \item Use retrieval-augmented approaches (RAG) for large corpora
          \item Summarize or chunk when possible
        \end{itemize}
      \end{alertblock}
    \end{column}
  \end{columns}
\end{frame}

% ---- Demo placeholder -----------------------------------------------
\demoframe{Iterative paper explanation \& hypothesis brainstorming — live with Claude}

% ============================================================
%  Section 4: Coding Assistants in Research
% ============================================================

\section{Coding Assistants in Research}

% ---- Why it matters ----------------------------------------
\begin{frame}[shrink=10]{Why Coding Assistance Matters for Researchers}
  \begin{columns}[T, onlytextwidth]
    \begin{column}{0.48\textwidth}
      \textbf{Researcher reality}\\[0.3em]
      \begin{itemize}
        \item Most CS/AI researchers write code daily
        \item Codebases are often solo or small-team maintained
        \item Prototype speed $\Rightarrow$ experiment iteration rate
        \item Debugging detracts from research thinking
        \item Large inherited codebases are hard to onboard
      \end{itemize}
    \end{column}
    \hfill
    \begin{column}{0.48\textwidth}
      \textbf{What AI assistance changes}\\[0.3em]
      \begin{itemize}
        \item Baseline implementations in hours, not days
        \item Explain unfamiliar library APIs instantly
        \item Auto-generate boilerplate (data loaders, eval loops)
        \item Catch bugs via conversational debugging
        \item Translate pseudocode $\to$ runnable code
      \end{itemize}
      \vspace{0.5em}
      \parencite{dakhel2023copilot}: Copilot reduces time-to-completion but requires careful review of generated code.
    \end{column}
  \end{columns}
\end{frame}

% ---- Spotlight: GitHub Copilot ------------------------------
\begin{frame}[shrink=10]{Spotlight: \tool{GitHub Copilot}}
  \begin{columns}[T, onlytextwidth]
    \begin{column}{0.55\textwidth}
      \textbf{Core capabilities}\\[0.3em]
      \begin{itemize}
        \item Inline code completion (tab to accept)
        \item Copilot Chat: explain, fix, test generation
        \item PR summarization and code review
        \item Agents mode (multi-file edits)
        \item Available in VS Code, JetBrains, Vim/Neovim
      \end{itemize}
      \vspace{0.5em}
      \textbf{Research use cases}\\[0.3em]
      \begin{itemize}
        \item Fast boilerplate: PyTorch training loops, HuggingFace pipelines
        \item Unit test generation for research code
        \item Documentation of legacy research repos
      \end{itemize}
    \end{column}
    \hfill
    \begin{column}{0.42\textwidth}
      \textbf{Strengths}\\
      \begin{itemize}
        \item Tightly integrated; near-zero context switching
        \item Large training corpus — knows most popular libraries
        \item Enterprise tier with IP indemnification
      \end{itemize}
      \vspace{0.5em}
      \textbf{Limitations}\\
      \begin{itemize}
        \item Code sent to GitHub servers
        \item Weaker at novel algorithms (not in training data)
        \item Suggestions can be plausible but subtly wrong~\parencite{barke2023grounded}
      \end{itemize}
    \end{column}
  \end{columns}
\end{frame}

% ---- Spotlight: Cursor -------------------------------------
\begin{frame}{Spotlight: \tool{Cursor}}
  \begin{columns}[T, onlytextwidth]
    \begin{column}{0.55\textwidth}
      \textbf{Key features}\\[0.3em]
      \begin{itemize}
        \item \textbf{Composer} (Agent mode): multi-file edits from a single prompt
        \item \textbf{Codebase indexing}: semantic search over entire repo
        \item \textbf{@-mentions}: reference specific files, docs, web in prompts
        \item Supports multiple backends: Claude, GPT-4o, local
        \item VS Code fork — full extension compatibility
      \end{itemize}
      \vspace{0.4em}
      \textbf{Research use cases}\\[0.3em]
      \begin{itemize}
        \item ``Implement Algorithm 2 from the attached paper''
        \item Refactor a training script across 10 files
        \item Ask questions about a large inherited codebase
      \end{itemize}
    \end{column}
    \hfill
    \begin{column}{0.42\textwidth}
      \small
      \textbf{Strengths}\\
      \begin{itemize}
        \item Best-in-class multi-file editing UX
        \item Codebase context without manual copy-paste
        \item Rapid iteration for exploratory coding
      \end{itemize}
      \vspace{0.2em}
      \begin{alertblock}{Privacy note}
        Code is sent to Cursor's servers (and upstream providers). For sensitive IP or proprietary data, check your institution's policy or use Privacy Mode.
      \end{alertblock}
    \end{column}
  \end{columns}
\end{frame}

% ---- Spotlight: Claude Code --------------------------------
\begin{frame}[fragile, shrink=15]{Spotlight: \tool{Claude Code}}
  \begin{columns}[T, onlytextwidth]
    \begin{column}{0.55\textwidth}
      \textbf{Architecture}\\[0.3em]
      \begin{itemize}
        \item Terminal-native agentic coding assistant
        \item Full file read/write, shell command execution
        \item Git-aware: diffs, commits, branch management
        \item \textbf{CLAUDE.md}: project-level instructions and context
        \item Extended thinking: configurable reasoning budget
      \end{itemize}
      \vspace{0.4em}
      \textbf{Research use cases}\\[0.3em]
      \begin{itemize}
        \item ``Implement this baseline end-to-end from the paper''
        \item Automated experiment script generation
        \item Data preprocessing pipeline construction
        \item Generate and run unit tests on research code
      \end{itemize}
    \end{column}
    \hfill
    \begin{column}{0.42\textwidth}
      \small
      \textbf{Distinctive features}\\[0.2em]
      \begin{itemize}
        \item Works in the terminal — no IDE lock-in
        \item Reads entire codebases before acting
        \item Interruptible mid-task
        \item IDE integrations: VS Code, JetBrains
      \end{itemize}
      \vspace{0.2em}
      \begin{lstlisting}[language=bash, numbers=none]
# Install
npm install -g @anthropic-ai/claude-code

# Run in your research repo
cd my-research-project
claude
      \end{lstlisting}
    \end{column}
  \end{columns}
\end{frame}

% ---- Spotlight: Cody / Aider -------------------------------
\begin{frame}[fragile]{Spotlight: \tool{Cody} \& \tool{Aider}}
  \begin{columns}[T, onlytextwidth]
    \begin{column}{0.48\textwidth}
      \begin{block}{\tool{Cody} (Sourcegraph)}
        \begin{itemize}
          \item Open-source, model-agnostic
          \item Indexes large codebases (millions of lines)
          \item \textbf{Best for:} navigating large inherited research frameworks (e.g., Fairseq, Megatron-LM)
          \item Supports Claude, GPT-4o, local models
          \item IDE plugins: VS Code, JetBrains
        \end{itemize}
      \end{block}
    \end{column}
    \hfill
    \begin{column}{0.48\textwidth}
      \begin{block}{\tool{Aider}}
        \begin{itemize}
          \item Git-native CLI coding assistant
          \item Maps entire repo before editing
          \item Commits changes automatically with meaningful messages
          \item Supports 100+ LLMs via litellm
          \item \textbf{Best for:} power users comfortable in the terminal
        \end{itemize}
        \begin{lstlisting}[language=bash, numbers=none]
pip install aider-chat
aider --model claude-3-5-sonnet
        \end{lstlisting}
      \end{block}
    \end{column}
  \end{columns}
\end{frame}

% ---- Comparison matrix -------------------------------------
\begin{frame}{Coding Assistant Comparison}
  \small
  \begin{tabularx}{\textwidth}{lYYYY}
    \toprule
    \textbf{Tool} & \textbf{Agentic} & \textbf{Open source} & \textbf{Multi-file} & \textbf{Best for researchers} \\
    \midrule
    \tool{Copilot}      & Partial & No  & Partial & Fast inline completion \\
    \tool{Cursor}       & Yes     & No  & Yes     & GUI-based multi-file edits \\
    \tool{Claude Code}  & Yes     & No  & Yes     & Terminal, end-to-end tasks \\
    \tool{Cody}         & Partial & Yes & Yes     & Large inherited codebases \\
    \tool{Aider}        & Yes     & Yes & Yes     & Git-native CLI power users \\
    \tool{Continue}     & Partial & Yes & Partial & Model-agnostic IDE plugin \\
    \bottomrule
  \end{tabularx}
  \vspace{0.5em}
  \small \alert{No single winner} — the right tool depends on your workflow, IDE preference, and data-sensitivity requirements.
\end{frame}

% ---- Demo placeholder -----------------------------------------------
\demoframe{Claude Code — implementing a research baseline from scratch, end-to-end}

% ============================================================
%  Section 5: Agentic Workflows
% ============================================================

\section{Agentic Workflows}

% ---- What is an agent? ------------------------------------
\begin{frame}{What Is an Agent?}
  \begin{columns}[T, onlytextwidth]
    \begin{column}{0.50\textwidth}
      \textbf{Definition}\\[0.3em]
      An \textbf{agent} is an LLM that can:
      \begin{enumerate}
        \item \textbf{Observe} its environment (context, tool outputs)
        \item \textbf{Plan} a sequence of actions
        \item \textbf{Act} by invoking tools
        \item \textbf{Iterate} until the goal is reached
      \end{enumerate}
      \vspace{0.5em}
      This is the \textbf{ReAct} framework~\parencite{yao2023react}: interleaving \emph{reasoning} and \emph{action}.
    \end{column}
    \hfill
    \begin{column}{0.46\textwidth}
      \centering
      \begin{tikzpicture}[
        node distance=0.9cm,
        box/.style={draw, rounded corners=3pt, fill=mLightBrown!15,
                    minimum width=2.0cm, minimum height=0.7cm,
                    align=center, font=\small},
        arr/.style={-{Stealth[length=4pt]}, thick, mLightBrown},
      ]
        \node[box] (obs)  {Observe};
        \node[box, below=of obs]  (plan) {Plan};
        \node[box, below=of plan] (act)  {Act};
        \node[box, below=of act]  (tool) {Tools\\(search, code, file)};
        \node[box, below=of tool] (eval) {Evaluate\\output};

        \draw[arr] (obs.south)  -- (plan.north);
        \draw[arr] (plan.south) -- (act.north);
        \draw[arr] (act.south)  -- (tool.north);
        \draw[arr] (tool.south) -- (eval.north);
        \draw[arr] (eval.east)  -- ++(0.8,0) |- (obs.east)
          node[pos=0.15, right, font=\tiny]{loop};
      \end{tikzpicture}
    \end{column}
  \end{columns}
\end{frame}

% ---- Tool use — building block ----------------------------
\begin{frame}[fragile]{Tool Use — The Building Block}
  \begin{columns}[T, onlytextwidth]
    \begin{column}{0.55\textwidth}
      \textbf{Mechanism}\\[0.3em]
      \begin{itemize}
        \item Model receives a list of available tools as JSON schemas
        \item Model outputs a structured tool-call instead of plain text
        \item Host executes the call, returns results to the model
        \item Model continues reasoning with the result~\parencite{schick2023toolformer}
      \end{itemize}
      \vspace{0.5em}
      \textbf{Common research tools}\\[0.3em]
      \begin{itemize}
        \item \texttt{search\_arxiv(query)} — fetch paper metadata
        \item \texttt{execute\_python(code)} — run analysis code
        \item \texttt{read\_file(path)} / \texttt{write\_file(path, content)}
        \item \texttt{run\_shell(cmd)} — arbitrary commands
        \item \texttt{mcp\_server} — connect to external tools
        \item \textbf{Google Antigravity} — agentic dev platform (IDE)
      \end{itemize}
    \end{column}
    \hfill
    \begin{column}{0.42\textwidth}
      \textbf{Tool schema example}\\[0.3em]
      \begin{lstlisting}[language=json, numbers=none, basicstyle=\ttfamily\scriptsize]
{
  "name": "search_arxiv",
  "description": "Search arXiv for papers",
  "input_schema": {
    "type": "object",
    "properties": {
      "query": {
        "type": "string",
        "description": "Search query"
      },
      "max_results": {
        "type": "integer",
        "default": 5
      }
    },
    "required": ["query"]
  }
}
      \end{lstlisting}
    \end{column}
  \end{columns}
\end{frame}

% ---- Model Context Protocol --------------------------------
\begin{frame}[shrink=18]{Model Context Protocol (MCP)}
  \begin{columns}[T, onlytextwidth]
    \begin{column}{0.55\textwidth}
      \textbf{What is MCP?}~\parencite{anthropic2024mcp}\\[0.3em]
      \begin{itemize}
        \item Open standard for LLM $\leftrightarrow$ external resource connectivity
        \item Proposed by Anthropic; adopted by \tool{Cursor}, \tool{Continue}, \tool{Claude Code}, VS Code, JetBrains
        \item Decouples the \emph{host} (LLM application) from \emph{servers} (data/tool providers)
        \item Communication via JSON-RPC 2.0 over stdio or SSE
        \item No widely-adopted competing open standard; OpenAI's tools API is proprietary
      \end{itemize}
      \vspace{0.5em}
      \textbf{Servers expose three primitives}\\[0.3em]
      \begin{itemize}
        \item \textbf{Resources}: data the model can read (files, DBs)
        \item \textbf{Tools}: functions the model can call
        \item \textbf{Prompts}: reusable prompt templates
      \end{itemize}
    \end{column}
    \hfill
    \begin{column}{0.42\textwidth}
      \centering
      \begin{tikzpicture}[
        node distance=0.8cm and 1.2cm,
        box/.style={draw, rounded corners=2pt, minimum width=2.2cm,
                    minimum height=0.65cm, align=center, font=\tiny},
        arr/.style={-{Stealth[length=4pt]}, thick},
      ]
        \node[box, fill=mLightBrown!20] (host)   {LLM Host\\(Claude Code)};
        \node[box, fill=blue!15, below=of host]   (client) {MCP Client};
        \node[box, fill=green!15, below=of client](server) {MCP Server};
        \node[box, fill=gray!15, below left=0.5cm and -0.4cm of server]  (r1) {Resources};
        \node[box, fill=gray!15, below=0.5cm of server]                  (r2) {Tools};
        \node[box, fill=gray!15, below right=0.5cm and -0.4cm of server] (r3) {Prompts};

        \draw[arr, <->] (host.south)   -- (client.north) node[midway,right,font=\tiny]{};
        \draw[arr, <->] (client.south) -- (server.north) node[midway,right,font=\tiny]{JSON-RPC};
        \draw[arr] (server.south west) -- (r1.north);
        \draw[arr] (server.south)      -- (r2.north);
        \draw[arr] (server.south east) -- (r3.north);
      \end{tikzpicture}
    \end{column}
  \end{columns}
\end{frame}

% ---- Research applications ---------------------------------
\begin{frame}{Research Applications of Agents}
  \begin{columns}[T, onlytextwidth]
    \begin{column}{0.48\textwidth}
      \begin{block}{Literature Sweeps}
        Agent queries arXiv/Semantic Scholar, filters by relevance score, downloads PDFs, extracts structured information, and produces a synthesis report.
      \end{block}
      \vspace{0.4em}
      \begin{block}{Experiment Orchestration}
        Agent reads experiment config, modifies hyperparameters, launches training runs (via \texttt{subprocess}), monitors logs, and emails results.
      \end{block}
    \end{column}
    \hfill
    \begin{column}{0.48\textwidth}
      \begin{block}{Data Pipeline}
        Agent writes a preprocessing script, validates outputs, generates summary statistics, and flags anomalies for human review.
      \end{block}
      \vspace{0.4em}
      \begin{block}{Report Generation}
        Agent reads experiment results, produces tables and figures (matplotlib), and writes a draft results section in LaTeX.
      \end{block}
    \end{column}
  \end{columns}
\end{frame}

% ---- Multi-agent patterns ----------------------------------
\begin{frame}{Multi-Agent Patterns}
  \begin{columns}[T, onlytextwidth]
    \begin{column}{0.50\textwidth}
      \textbf{Orchestrator + Subagents}\\[0.3em]
      \begin{tikzpicture}[
        node distance=0.6cm and 1.2cm,
        box/.style={draw, rounded corners=2pt, minimum width=1.8cm,
                    minimum height=0.6cm, align=center, font=\tiny,
                    fill=mLightBrown!15},
        arr/.style={-{Stealth[length=4pt]}, thick},
      ]
        \node[box] (orch) {Orchestrator LLM};
        \node[box, below left=0.6cm and 0.3cm of orch]  (s1) {Search\\Agent};
        \node[box, below=0.6cm of orch]                  (s2) {Code\\Agent};
        \node[box, below right=0.6cm and 0.3cm of orch] (s3) {Write\\Agent};
        \draw[arr] (orch.south west) -- (s1.north);
        \draw[arr] (orch.south)      -- (s2.north);
        \draw[arr] (orch.south east) -- (s3.north);
        \draw[arr, dashed] (s1.north) -- (orch.south west);
        \draw[arr, dashed] (s2.north) -- (orch.south);
        \draw[arr, dashed] (s3.north) -- (orch.south east);
      \end{tikzpicture}
    \end{column}
    \hfill
    \begin{column}{0.46\textwidth}
      \textbf{Parallel Subagents}\\[0.3em]
      \begin{itemize}
        \item Independent tasks run concurrently
        \item Orchestrator merges results
        \item Reduces total wall-clock time
        \item Example: 10 papers analyzed in parallel
        \item Frameworks: \tool{LangChain}, \tool{AutoGen}, \tool{Google Antigravity}
      \end{itemize}
      \vspace{0.5em}
      \textbf{Reference}\\
      \parencite{wang2024survey}: survey of LLM-based autonomous agents and multi-agent architectures.
    \end{column}
  \end{columns}
\end{frame}

% ---- Harness Engineering Slide ---------------------------
\begin{frame}{Ecosystem Shift: Prompt $\to$ Context $\to$ Harness Engineering}
  \begin{columns}[T, onlytextwidth]
    \begin{column}{0.48\textwidth}
      \textbf{The Harness as Core Object}\\
      The model is only one part of an agent. The \textbf{harness} controls:
      \begin{itemize}
        \item Tool access \& execution permissions
        \item Short-term \& summary memory
        \item Context compaction / token control
        \item Evaluation \& telemetry logging
        \item Multi-agent handoffs / orchestration
      \end{itemize}
      \vspace{0.3em}
      \textbf{Loop Engineering}:
      \begin{itemize}
        \item Run repeated build-test-fix loops autonomously
        \item Automated benchmark execution loops
        \item Specialized peer-reviewer loops
      \end{itemize}
    \end{column}
    \hfill
    \begin{column}{0.48\textwidth}
      \textbf{Foundational Cases / Workflows}
      \begin{itemize}
        \item \textbf{OpenAI's Codex CLI Technical Guidance}: Describes the harness as the agent loop orchestrating user, model, and tools.
        \item \textbf{LangGraph}: Positions itself as an orchestration runtime for streaming, multi-agent coordination, and persistence.
      \end{itemize}
      \vspace{0.4em}
      \begin{block}{Harness over Model}
        \small For research workflows, the harness matters as much as the model: it determines what gets remembered, what tools are safe to call, what gets retried, and what evidence is logged.
      \end{block}
    \end{column}
  \end{columns}
\end{frame}

% ---- Popular agent harnesses ------------------------------
\begin{frame}{Popular Agent Harnesses (2026)}
  \scriptsize
  \begin{tabularx}{\textwidth}{>{\raggedright\arraybackslash}p{0.18\textwidth} >{\raggedright\arraybackslash}p{0.32\textwidth} X}
    \toprule
    \textbf{Harness} & \textbf{Strength} & \textbf{Good fit for research} \\
    \midrule
    \tool{LangGraph} (LC) & Durable graph, persistence, Deep Agents & Complex pipelines with checkpoints, retries, and human approval gates. \\
    \tool{Pi} & Minimal, hackable; packageable skills/extensions & Researchers who want to customize/inspect the loop over using heavy frameworks. \\
    \tool{smolagents} (HF) & Minimal HF library, executable Python code & Quick prototypes and tool loops in a few lines of code. \\
    \tool{OpenAI Agents SDK} & Tool-first hosted agents & Strong when your stack is already centered on OpenAI tools and tracing. \\
    \tool{CrewAI} / AutoGen & Multi-agent conversation/role patterns & Lightweight team-style or role-collaborative (planner, coder, critic) workflows. \\
    \bottomrule
  \end{tabularx}
  \vspace{0.4em}
  \normalsize
  \textbf{Selection heuristic:} pick by \emph{control model} (graph vs minimal loop), \emph{observability}, \emph{human-in-the-loop support}, and \emph{customizability}.
\end{frame}

% ---- Agent Memory Slide ----------------------------------
\begin{frame}{Agent Memory Is Not Just ``Bigger Context''}
  \begin{columns}[T, onlytextwidth]
    \begin{column}{0.48\textwidth}
      \textbf{Three Memory Patterns}\\[0.3em]
      \begin{itemize}
        \item \textbf{Context Compaction}: Dynamically summarizing old interactions or compressing raw text to fit the attention window.
        \item \textbf{Project Memory}: File-based repo context, such as \texttt{CLAUDE.md}, repository maps, local notes, and current task state.
        \item \textbf{Structured Memory Graphs}: e.g., \tool{Beads} (VirtusLab), which provides persistent, structured issue-tracking memory and replaces flat markdown plans with a dependency-aware graph across model restarts.
      \end{itemize}
    \end{column}
    \hfill
    \begin{column}{0.48\textwidth}
      \textbf{Portable Agent Memory}\\[0.3em]
      Citing emerging research proposing structured memory transfer across heterogeneous agents with provenance and fine-grained access control.
      \vspace{0.4em}
      \begin{alertblock}{Warning: The Memory Trade-off}
        Memory increases continuity, but also increases \textbf{privacy risk, stale code assumptions, prompt-injection persistence}, and the \textbf{accidental reuse of wrong decisions}.
      \end{alertblock}
    \end{column}
  \end{columns}
\end{frame}

% ---- Persistent Agents Slide -----------------------------
\begin{frame}{Persistent Agents: Beyond the IDE}
  \small
  \begin{tabularx}{\textwidth}{p{0.22\textwidth} X}
    \toprule
    \textbf{System} & \textbf{Capabilities \& Focus} \\
    \midrule
    \tool{OpenClaw} & Self-hosted personal agent controlled from chat apps; manages email, calendar, files, and tools. Strong data-sovereignty story, but high security risk if over-permissioned (clears inboxes, sends mail). \\
    \midrule
    \tool{Hermes Agent} (Nous Research) & Self-hosted agent with persistent memory and skill creation. An agent that ``grows with you,'' creates new skills from experience, and searches past conversations. \\
    \bottomrule
  \end{tabularx}
  \vspace{0.6em}
  \begin{alertblock}{Key Practice Lesson}
    Persistent operational agents require strict \textbf{permission boundaries, detailed audit logs, memory hygiene}, and \textbf{failsafe revocation controls} to prevent runaway processes or security breaches.
  \end{alertblock}
\end{frame}

% ---- Current limitations -----------------------------------
\begin{frame}{Current Limitations}
  \begin{columns}[T, onlytextwidth]
    \begin{column}{0.48\textwidth}
      \begin{alertblock}{Reliability issues}
        \begin{itemize}
          \item \textbf{Error propagation}: mistakes in early steps compound; agents rarely self-correct without human intervention
          \item \textbf{Cost}: multi-step agents consume many tokens; long tasks can be expensive
        \end{itemize}
      \end{alertblock}
    \end{column}
    \hfill
    \begin{column}{0.48\textwidth}
      \begin{alertblock}{Verification gaps}
        \begin{itemize}
          \item \textbf{Verification}: agents cannot always verify their own outputs; scientific claims must be checked
          \item \textbf{Human-in-the-loop}: for consequential actions (delete files, push to git, email), always require explicit approval
        \end{itemize}
      \end{alertblock}
    \end{column}
  \end{columns}
  \vspace{0.6em}
  \textbf{Best practice:} Use agents for \emph{reversible, low-stakes} tasks. Build in checkpoints for irreversible actions.
  \vspace{0.4em}

  \textbf{Progress note (early 2026):} Coding agents have improved dramatically. SWE-agent v1.07 resolves $\sim$68\% of SWE-bench Verified issues; Claude Opus 4.8 (Thinking) reaches 83.1\%; GPT-5.5 achieves 79.5\% on Terminal-Bench 2.0 and 59.4\% on SWE-Bench Pro. Human oversight remains essential for production code. \textit{(As of June 2026)}
\end{frame}

% ---- Demo idea -------------------------------------------------------
\begin{frame}[shrink=8]{Live Demo Idea: Research Harness on Planning Benchmarks}
  \textbf{Goal (7--8 min):} show one harness coordinating search, coding, and writing in a single loop.

  \vspace{0.4em}
  \textbf{Scenario (from this repo)}
  \begin{itemize}
    \item Start in the Pyperplan workspace (Section 5 context)
    \item Task: compare 2 heuristics on one benchmark domain and produce a mini result summary
  \end{itemize}

  \vspace{0.3em}
  \textbf{Harness roles}
  \begin{itemize}
    \item \textbf{Planner agent}: chooses benchmark/problem pairs and run commands
    \item \textbf{Coder agent}: patches experiment script + parses run logs into a small CSV
    \item \textbf{Writer agent}: drafts 5 bullet points + 1 LaTeX-ready table row
  \end{itemize}

  \vspace{0.3em}
  \textbf{Why it connects to this talk}
  \begin{itemize}
    \item \textbf{Coding assistants}: concrete repo edits and terminal execution
    \item \textbf{Agentic workflows}: orchestrator + subagents + tool calls
    \item \textbf{Writing support}: automatic conversion of results into slide/paper text
  \end{itemize}

  \vspace{0.3em}
  \textbf{Safety gate:} require manual approval before any file overwrite outside a demo folder.
\end{frame}

% ============================================================
%  Section 6: Literature and Research Discovery
% ============================================================

\section{Literature \& Research Discovery}

% ---- The discovery problem ---------------------------------
\begin{frame}{The Discovery Problem}
  \begin{columns}[T, onlytextwidth]
    \begin{column}{0.55\textwidth}
      \textbf{Scale of the challenge}\\[0.3em]
      \begin{itemize}
        \item arXiv: $>$2.97 million papers as of early 2026; $\sim$28K new submissions/month (all categories)
        \item PubMed: $>$39 million citations
        \item Keyword search misses semantically related work
        \item Citation networks reveal influence but not recency
        \item Staying current requires curated workflows
      \end{itemize}
      \vspace{0.5em}
      \textbf{Traditional approaches failing}\\[0.3em]
      \begin{itemize}
        \item Google Scholar: good recall, poor precision
        \item Manual RSS / email alerts: unscalable
        \item Conference proceedings: siloed by venue
      \end{itemize}
    \end{column}
    \hfill
    \begin{column}{0.42\textwidth}
      \textbf{Semantic vs.\ keyword search}\\[0.4em]
      \begin{tabular}{lcc}
        \toprule
        & \textbf{Keyword} & \textbf{Semantic} \\
        \midrule
        Exact matches & $\checkmark\checkmark$ & $\checkmark$ \\
        Related concepts & $\times$ & $\checkmark\checkmark$ \\
        Multilingual & $\times$ & $\checkmark$ \\
        Speed & Fast & Moderate \\
        Explainability & High & Lower \\
        \bottomrule
      \end{tabular}
    \end{column}
  \end{columns}
\end{frame}

% ---- Tool survey -------------------------------------------
\begin{frame}{Tool Survey: Research Discovery}
  \small
  \begin{tabularx}{\textwidth}{lXX}
    \toprule
    \textbf{Tool} & \textbf{Core strength} & \textbf{Best use case} \\
    \midrule
    \tool{Semantic Scholar} & Semantic search, citation graph, open API & Programmatic access, citation analysis \\
    \tool{Elicit}           & Structured extraction (PICO / columns) & Systematic reviews, hypothesis scoping \\
    \tool{ResearchRabbit}   & Citation network visualization, alerts & Snowballing from seed papers \\
    \tool{Consensus}        & Evidence synthesis, yes/no queries & Quick claim verification \\
    \tool{Connected Papers} & Visual citation graph & Discovering neighboring papers \\
    \tool{Perplexity}       & Web + academic search with citations & Fast question answering with sources \\
    \tool{NotebookLM}       & Document Q\&A with grounding & Deep dive into a curated paper set \\
    \bottomrule
  \end{tabularx}
\end{frame}

% ---- Spotlight: Elicit -------------------------------------
\begin{frame}{Spotlight: \tool{Elicit}}
  \begin{columns}[T, onlytextwidth]
    \begin{column}{0.55\textwidth}
      \small
      \textbf{What Elicit does}\\[0.2em]
      \begin{itemize}
        \item AI-powered systematic review assistant
        \item Input: research question $\to$ Elicit queries literature
        \item Extracts structured data: population, intervention,\\ outcome, sample size, findings
        \item Presents results as a \emph{configurable extraction table}
        \item Upload your own PDFs for extraction
      \end{itemize}
      \vspace{0.2em}
      \textbf{Researcher workflow}\\[0.2em]
      \begin{enumerate}
        \item Pose a specific research question
        \item Review top-N papers in table form
        \item Export to CSV for further analysis
        \item Synthesize patterns across studies
      \end{enumerate}
    \end{column}
    \hfill
    \begin{column}{0.42\textwidth}
      \textbf{Example extraction columns}\\[0.4em]
      \begin{tabular}{l}
        \toprule
        Paper title \& year \\
        Dataset used \\
        Method category \\
        Key metric \& value \\
        Reported limitations \\
        \bottomrule
      \end{tabular}
      \vspace{0.7em}
      \alert{Caveat:} Elicit only covers papers indexed by Semantic Scholar. Verify extractions — LLM-based extraction can miss nuance.\\
      \vspace{0.3em}
      \textbf{Pricing:} Free Basic plan; Plus (\$120/yr) and Pro (\$499/yr) plans for more projects, PDFs, and systematic review features.
    \end{column}
  \end{columns}
\end{frame}

% ---- Spotlight: ResearchRabbit ----------------------------
\begin{frame}{Spotlight: \tool{ResearchRabbit}}
  \begin{columns}[T, onlytextwidth]
    \begin{column}{0.55\textwidth}
      \small
      \textbf{Core feature: Citation Network Visualization}\\[0.2em]
      \begin{itemize}
        \item Start with 1--3 seed papers
        \item Visualizes papers citing and cited-by the seeds
        \item Color-codes by recency, citation count, and similarity
        \item ``This paper only'' vs.\ ``All similar work'' views
        \item \textbf{Zotero integration}: sync collections bi-directionally
      \end{itemize}
      \vspace{0.2em}
      \textbf{Use for snowball literature search}\\[0.2em]
      \begin{enumerate}
        \item Seed with 2--3 highly cited papers in your area
        \item Explore ``later work'' graph
        \item Filter by year, venue, or author
        \item Export selected papers to Zotero
      \end{enumerate}
    \end{column}
    \hfill
    \begin{column}{0.42\textwidth}
      \textbf{Strengths}\\[0.3em]
      \begin{itemize}
        \item Visual — reveals clusters of related work instantly
        \item Weekly email digests of new papers in your collection's neighborhood
        \item Freemium model (Nov 2025): free tier limited to 50 inputs / 1 project;\\
              ResearchRabbit+ (\$120/yr) for 300 inputs and multiple projects
      \end{itemize}
      \vspace{0.5em}
      \textbf{Limitations}\\[0.3em]
      \begin{itemize}
        \item Coverage dependent on Semantic Scholar
        \item Visualization can be overwhelming for large graphs
        \item No structured extraction (complement with Elicit)
      \end{itemize}
    \end{column}
  \end{columns}
\end{frame}

% ---- Spotlight: Semantic Scholar API -----------------------
\begin{frame}[fragile]{Spotlight: \tool{Semantic Scholar} + API}
  \begin{columns}[T, onlytextwidth]
    \begin{column}{0.55\textwidth}
      \textbf{Programmatic access}~\parencite{lo2023semantic}\\[0.3em]
      \begin{lstlisting}[language=python, basicstyle=\ttfamily\scriptsize]
import requests

BASE = "https://api.semanticscholar.org"

def search_papers(query, limit=10):
    url = BASE + "/graph/v1/paper/search"
    params = {
        "query": query,
        "limit": limit,
        "fields": "title,year,authors,"
                  "citationCount,abstract,"
                  "openAccessPdf"
    }
    r = requests.get(url, params=params)
    return r.json()["data"]

papers = search_papers(
    "transformer attention mechanism"
)
for p in papers:
    print(p["title"], p["year"],
          p["citationCount"])
      \end{lstlisting}
    \end{column}
    \hfill
    \begin{column}{0.42\textwidth}
      \textbf{API capabilities}\\[0.3em]
      \begin{itemize}
        \item Search by keyword, author, venue, DOI
        \item Fetch full citation/reference lists
        \item Paper recommendations via \texttt{/recommendations}
        \item Rate limit: 1 req/sec per API key (free); shared unauthenticated pool
        \item Open Access PDF URLs included when available
      \end{itemize}
      \vspace{0.4em}
      \textbf{Combine with LLMs:}\\
      Feed paper metadata to Claude for\\ relevance scoring and synthesis.
    \end{column}
  \end{columns}
\end{frame}

% ---- Workflow pipeline -------------------------------------
\begin{frame}{Research Discovery Workflow}
  \centering
  \resizebox{\textwidth}{!}{%
  \begin{tikzpicture}[
    node distance=0.5cm and 1.0cm,
    box/.style={draw, rounded corners=3pt, fill=mLightBrown!15,
                minimum width=2.2cm, minimum height=0.75cm,
                align=center, font=\small},
    arr/.style={-{Stealth[length=5pt]}, thick, mLightBrown},
    lbl/.style={font=\tiny, align=center},
  ]
    \node[box] (kw)   {Keyword\\search};
    \node[box, right=of kw]   (ss)  {Semantic\\Scholar};
    \node[box, right=of ss]   (cg)  {Citation\\graph\\(ResearchRabbit)};
    \node[box, right=of cg]   (ex)  {Structured\\extraction\\(Elicit)};
    \node[box, right=of ex]   (syn) {Synthesis\\(NotebookLM)};
    \node[box, right=of syn]  (val) {Manual\\validation};

    \draw[arr] (kw.east)  -- (ss.west);
    \draw[arr] (ss.east)  -- (cg.west);
    \draw[arr] (cg.east)  -- (ex.west);
    \draw[arr] (ex.east)  -- (syn.west);
    \draw[arr] (syn.east) -- (val.west);
  \end{tikzpicture}}%

  \vspace{0.4em}
  \begin{itemize}
    \item The pipeline is \textbf{iterative} — new papers found in step 3 seed another round
    \item \textbf{Human validation is non-negotiable} at every stage
    \item AI tools speed up steps 1--5; step 6 remains your responsibility
  \end{itemize}
\end{frame}

% ---- Demo placeholder -----------------------------------------------
\demoframe{Elicit structured review + ResearchRabbit citation network — live demonstration}

% ============================================================
%  Section 7: Writing Assistance
% ============================================================

\section{Writing Assistance}

% ---- Where AI writing help is legitimate -------------------
\begin{frame}{Where AI Writing Help Is Legitimate}
  \begin{columns}[T, onlytextwidth]
    \begin{column}{0.48\textwidth}
      \textbf{Pre-writing}\\
      \begin{itemize}
        \item Outline generation from notes
        \item Brainstorming the narrative arc of a paper
        \item Identifying gaps in an argument structure
        \item Generating alternative framings of a contribution
      \end{itemize}
      \vspace{0.5em}
      \textbf{Drafting assistance}\\
      \begin{itemize}
        \item Expand bullet points into prose
        \item Draft related work from citations
        \item Write boilerplate sections (data collection, ethics)
      \end{itemize}
    \end{column}
    \hfill
    \begin{column}{0.48\textwidth}
      \textbf{Editing \& revision}\\
      \begin{itemize}
        \item Grammar and clarity improvement
        \item Register adjustment (formal/informal)
        \item Shortening verbose passages
        \item Non-native speaker polishing
      \end{itemize}
      \vspace{0.5em}
      \textbf{Specific tasks}\\
      \begin{itemize}
        \item Abstract generation from full paper
        \item Title and subtitle brainstorming
        \item Response-to-reviewers drafts
        \item Grant proposal narratives
      \end{itemize}
    \end{column}
  \end{columns}
\end{frame}

% ---- Prompting patterns for academic writing ---------------
\begin{frame}[fragile]{Prompting Patterns for Academic Writing}
  \begin{columns}[T, onlytextwidth]
    \begin{column}{0.48\textwidth}
      \begin{lstlisting}[language={}, numbers=none, basicstyle=\ttfamily\scriptsize,
                         caption={Abstract generation}]
You are a technical writer for a top ML venue.
Here is the full paper:
[PAPER TEXT]

Write a 150-word abstract following the
structure: motivation, problem, method,
results, contribution. Use precise
terminology from the paper.
      \end{lstlisting}
      \vspace{0.3em}
      \begin{lstlisting}[language={}, numbers=none, basicstyle=\ttfamily\scriptsize,
                         caption={Related work expansion}]
Given these 5 citation entries:
[BIB ENTRIES]

Write 2 paragraphs of related work
situating our contribution. Note how
each prior work differs from ours.
Method: [METHOD DESCRIPTION]
      \end{lstlisting}
    \end{column}
    \hfill
    \begin{column}{0.48\textwidth}
      \begin{lstlisting}[language={}, numbers=none, basicstyle=\ttfamily\scriptsize,
                         caption={Reviewer response}]
Here is Reviewer 2's comment:
[COMMENT]

And here is our response draft:
[DRAFT]

As a skeptical reviewer, identify
any gaps in our response. Suggest
specific improvements.
      \end{lstlisting}
      \vspace{0.3em}
      \begin{lstlisting}[language={}, numbers=none, basicstyle=\ttfamily\scriptsize,
                         caption={Clarity editing}]
Revise the following paragraph for
clarity and concision. Preserve all
technical content. Target: 20% shorter.
Active voice where possible.

[PARAGRAPH TEXT]
      \end{lstlisting}
    \end{column}
  \end{columns}
\end{frame}

% ---- Tool comparison for writing ---------------------------
\begin{frame}{Tool Comparison for Academic Writing}
  \footnotesize
  \begin{tabularx}{\textwidth}{lYYYY}
    \toprule
    \textbf{Tool} & \textbf{Long-form drafting} & \textbf{Technical accuracy} & \textbf{Citations} & \textbf{Best for} \\
    \midrule
    \tool{Claude 4.6 Sonnet} & Excellent & High  & No (verify) & All-round academic writing \\
    \tool{GPT-5 / 5.3-Codex} & Excellent & High  & No (verify) & Writing + coding tasks; 400K context \\
    \tool{Gemini 3.0 / 3.1 Pro} & Good      & Good  & Google Scholar & Research + search; best for long-context docs \\
    \tool{Grammarly}        & Grammar only & N/A & No & Grammar, style checking \\
    \tool{Writefull}        & Good      & Academic-tuned & No & Non-native speaker editing \\
    \bottomrule
  \end{tabularx}
  \vspace{0.2em}
  \footnotesize \alert{Note:} None of these tools produce verifiable citations independently. Always cross-check any citation they suggest.
\end{frame}

% ---- LaTeX-specific assistance ----------------------------
\begin{frame}[fragile]{LaTeX-Specific Assistance}
  \begin{columns}[T, onlytextwidth]
    \begin{column}{0.48\textwidth}
      {\small
      \textbf{What LLMs do well in LaTeX}\\[0.2em]
      \begin{itemize}
        \item Equation formatting from handwritten / pseudocode notation
        \item TikZ diagram generation from descriptions
        \item BibTeX entry cleanup and standardization
        \item Table construction (booktabs from CSV data)
        \item Algorithm environment from pseudocode
      \end{itemize}
      \vspace{0.3em}
      \textbf{Example prompt}\\[0.2em]}
      \begin{lstlisting}[language={}, numbers=none, basicstyle=\ttfamily\scriptsize]
Write a TikZ diagram of a
Transformer encoder block
showing: multi-head attention,
layer norm, FFN, residual
connections. Use a clean style.
      \end{lstlisting}
    \end{column}
    \hfill
    \begin{column}{0.48\textwidth}
      {\small\textbf{BibTeX cleanup prompt}\\[0.2em]}
      \begin{lstlisting}[language={}, numbers=none, basicstyle=\ttfamily\scriptsize]
Clean and standardize these BibTeX
entries. Fix: missing fields, title
casing (keep proper nouns), author
name format (Last, First). Use
consistent citation keys: AuthorYearKeyword.

[BIB ENTRIES]
      \end{lstlisting}
      \vspace{0.3em}
      {\small\alert{Compile and verify all generated LaTeX.} LLMs produce syntactically plausible but sometimes broken code — always test in your document.}
    \end{column}
  \end{columns}
\end{frame}

% ---- What to watch out for ---------------------------------
\begin{frame}{What to Watch Out For}
  \begin{columns}[T, onlytextwidth]
    \begin{column}{0.32\textwidth}
      \begin{alertblock}{Style Homogenization}
        AI-assisted papers can converge on a ``house style'' — smooth but generic prose. Preserve your voice. Use AI to improve clarity, not replace personality.
      \end{alertblock}
    \end{column}
    \hfill
    \begin{column}{0.32\textwidth}
      \begin{alertblock}{Over-Reliance}
        Heavy AI use can atrophy scientific writing skills. The act of writing \textbf{is} thinking; outsourcing it entirely means outsourcing thought. Use as an assistant, not a ghostwriter.
      \end{alertblock}
    \end{column}
    \hfill
    \begin{column}{0.32\textwidth}
      \begin{alertblock}{Confidentiality}
        Unpublished results, proprietary data, and confidential reviewer comments must \textbf{not} be pasted into commercial AI tools without institutional approval. Check your institution's AI policy.
      \end{alertblock}
    \end{column}
  \end{columns}
\end{frame}

% ============================================================
%  Section 8: Caveats, Limitations, and Ethics
% ============================================================

\section{Caveats, Limitations \& Ethics}

% ---- Hallucination ------------------------------------------
\begin{frame}{Hallucination — The Fundamental Problem}
  \begin{columns}[T, onlytextwidth]
    \begin{column}{0.55\textwidth}
      {\small
      \textbf{What hallucination is}\\[0.2em]
      \begin{itemize}
        \item LLMs generate text that is \emph{fluent but false}~\parencite{huang2023hallucination}
        \item Fabricated citations with real-looking DOIs, authors, venues
        \item Plausible-sounding but incorrect mathematical statements
        \item Outdated information presented as current
        \item Confident tone regardless of accuracy
      \end{itemize}
      \vspace{0.3em}
      \textbf{Particularly dangerous in research}\\[0.2em]
      \begin{itemize}
        \item Fabricated citations erode scientific record
        \item Wrong experimental details can waste resources
        \item False ``prior work'' shapes research direction incorrectly
      \end{itemize}}
    \end{column}
    \hfill
    \begin{column}{0.42\textwidth}
      \begin{block}{Mitigation strategies}
        {\small\begin{itemize}
          \item \textbf{Verify every citation} independently via DOI / Scholar
          \item Use models with \textbf{retrieval grounding} (Perplexity, NotebookLM)
          \item Ask model to \textbf{express uncertainty}: ``If you are not sure, say so''
          \item \textbf{Cross-check} key claims with a second model
          \item Prefer grounded tasks (summarize \emph{this} text) over generative tasks (tell me about X)
        \end{itemize}}
      \end{block}
    \end{column}
  \end{columns}
\end{frame}

% ---- Reproducibility ----------------------------------------
\begin{frame}[fragile]{Reproducibility}
  \begin{columns}[T, onlytextwidth]
    \begin{column}{0.55\textwidth}
      {\small
      \textbf{The problem}\\[0.2em]
      \begin{itemize}
        \item Prompts + model version determine outputs
        \item Models are updated silently — same prompt, different answer
        \item Non-determinism even at temperature=0 (infrastructure variance)
        \item AI-generated code may not reproduce without the exact model
      \end{itemize}
      \vspace{0.3em}
      \textbf{Best practices}~\parencite{pineau2021reproducibility}\\[0.2em]
      \begin{itemize}
        \item \textbf{Pin model versions} in your methods: \model{claude-fable-5}
        \item \textbf{Log all prompts} used in the research process
        \item \textbf{Declare AI use} in acknowledgments / methods
        \item Archive prompts in supplementary materials
      \end{itemize}}
    \end{column}
    \hfill
    \begin{column}{0.42\textwidth}
      \begin{block}{Methods section template}
        \small \textit{``We used [Model, Version, Date] to assist with [specific task]. The prompts are provided in Appendix~X. All AI-generated content was verified by the authors.''}
      \end{block}
      \vspace{0.3em}
      \textbf{Logging example}\\[0.2em]
      \begin{lstlisting}[language=python, numbers=none, basicstyle=\ttfamily\scriptsize]
import json, datetime
log = {
  "model": "claude-fable-5",
  "date": str(datetime.date.today()),
  "task": "related work drafting",
  "prompt": prompt_text,
}
json.dump(log, open("ai_log.json","a"))
      \end{lstlisting}
    \end{column}
  \end{columns}
\end{frame}

% ---- IP and data privacy -----------------------------------
\begin{frame}{IP and Data Privacy}
  \begin{columns}[T, onlytextwidth]
    \begin{column}{0.48\textwidth}
      {\small
      \textbf{Training data provenance}\\[0.2em]
      \begin{itemize}
        \item Frontier models trained on large web corpora; copyright status contested
        \item Generated code may resemble training data (memorization)
        \item IP ownership of AI-generated content is legally unsettled
      \end{itemize}
      \vspace{0.3em}
      \textbf{Terms of service}\\[0.2em]
      \begin{itemize}
        \item Commercial tools (GPT-5.6, Claude Fable 5) may use inputs for model improvement (check ToS)
        \item Enterprise / API tiers often have stronger privacy guarantees
        \item Institutional licenses may have different terms
      \end{itemize}}
    \end{column}
    \hfill
    \begin{column}{0.48\textwidth}
      \begin{block}{When to use local models}
        {\small Use open-weight models locally (\tool{Ollama}, \tool{vLLM}) for:
        \begin{itemize}
          \item Unpublished research data
          \item Confidential datasets (medical, industrial)
          \item Code under NDA
          \item Proprietary algorithms
          \item Pre-publication manuscripts
        \end{itemize}}
      \end{block}
      \vspace{0.3em}
      {\small\alert{Check your institution's AI policy before using any external tool with research data.}}
    \end{column}
  \end{columns}
\end{frame}

% ---- Peer review integrity ---------------------------------
\begin{frame}{Peer Review Integrity}
  \begin{columns}[T, onlytextwidth]
    \begin{column}{0.55\textwidth}
      \textbf{Venue policies (as of 2026)}\\[0.4em]
      {\footnotesize
      \begin{tabular}{@{}lp{4.8cm}@{}}
        \toprule
        \textbf{Venue} & \textbf{Policy} \\
        \midrule
        NeurIPS & Disclosure required (methods section) \\
        ICML    & Disclosure required (similar to NeurIPS) \\
        ICLR    & Disclosure required; prompt injection banned \\
        ACL     & Disclosure required; transparent use only \\
        Nature  & Disclose use; no AI co-authorship \\
        Science & AI-generated text prohibited; no AI author \\
        \bottomrule
      \end{tabular}}
    \end{column}
    \hfill
    \begin{column}{0.42\textwidth}
      \textbf{Key distinctions}\\[0.3em]
      \begin{block}{Permitted: assisting}
        Grammar checking, clarity editing, formatting, generating outlines you then rewrite.
      \end{block}
      \vspace{0.4em}
      \begin{alertblock}{Prohibited: delegating}
        Using AI to write reviews, fabricate data, generate conclusions you present as your own without verification.
      \end{alertblock}
      \vspace{0.4em}
      \small Norms are evolving. When in doubt, disclose.
    \end{column}
  \end{columns}
\end{frame}

% ---- Benchmark contamination --------------------------------
\begin{frame}{Benchmark Contamination \& Evaluation Validity}
  \begin{columns}[T, onlytextwidth]
    \begin{column}{0.55\textwidth}
      \textbf{The contamination problem}\\[0.3em]
      \begin{itemize}
        \item Training data includes public benchmarks (MMLU, HumanEval, etc.)
        \item Models may ``remember'' answers rather than solve them
        \item Performance inflation: reported scores may not reflect true capability
        \item Dynamic benchmarks (e.g., LiveCodeBench) partially mitigate this
      \end{itemize}
      \vspace{0.5em}
      \textbf{For your own evaluations}\\[0.3em]
      \begin{itemize}
        \item Use \textbf{held-out, private test sets} for LLM-based pipelines
        \item Do not use LLM-generated text as ground truth for LLM evaluation (circular)
        \item Consider task-specific human evaluation for qualitative claims
      \end{itemize}
    \end{column}
    \hfill
    \begin{column}{0.42\textwidth}
      \begin{alertblock}{Evaluating LLM-based systems}
        \begin{itemize}
          \item \textbf{Model-as-judge} has systematic biases (verbosity, self-preference)
          \item Always supplement with human evaluation~\parencite{liang2022holistic}
          \item Report inter-rater agreement
          \item Use blind evaluation where possible
        \end{itemize}
      \end{alertblock}
    \end{column}
  \end{columns}
\end{frame}

% ---- Agentic Security Slide ------------------------------
\begin{frame}[shrink=15]{Security in Agentic Systems: Critical Risks}
  \begin{columns}[T, onlytextwidth]
    \begin{column}{0.48\textwidth}
      \textbf{Key Threat Verticals}
      \begin{itemize}\setlength{\itemsep}{2pt}
        \item \textbf{Indirect Prompt Injection}: Malicious inputs in raw sources (emails, papers, GitHub issues) override agent instructions, executing unauthorized commands.
        \item \textbf{Over-Permissioned Tools}: Providing write, delete, or executables (e.g. bash) access without constraints is highly risky.
        \item \textbf{Credential Leaks}: Agent logs or context dumps accidentally printing API keys, environment variables, or private SSH keys.
        \item \textbf{Trust the Harness}: Trust the client and MCP tools you grant access. A \href{https://thereallo.dev/blog/claude-code-prompt-steganography\#marker}{2026 Claude Code analysis} reported subtle prompt markers tied to a custom API URL and timezone.
      \end{itemize}
    \end{column}
    \hfill
    \begin{column}{0.48\textwidth}
      \textbf{Supply Chain \& Action Risks}
      \begin{itemize}\setlength{\itemsep}{2pt}
        \item \textbf{Package Hallucination}: Agents requesting non-existent packages, enabling attackers to register them (dependency slopsquatting / typosquatting).
        \item \textbf{Irreversible Actions}: Sending emails, deleting local directories, database drops, or pulling untrusted scripts.
      \end{itemize}
      \vspace{0.4em}
      \begin{alertblock}{Best Mitigation Practice}
        Run agents in sandboxed containers (Docker/gVisor), require human-in-the-loop approvals for all write/external actions, clean system logs of any API files/secrets
      \end{alertblock}
    \end{column}
  \end{columns}
\end{frame}

% ---- Broader picture ----------------------------------------
\begin{frame}{The Broader Picture}
  \begin{columns}[T, onlytextwidth]
    \begin{column}{0.32\textwidth}
      \begin{block}{Compute \& Climate}
        Training costs remain large, but inference is now far cheaper. Recent estimates: $\sim$0.24--0.34\,Wh per text query ($\approx$0.03--1\,g CO\textsubscript{2}). Text generation per page $\approx$1\,g CO\textsubscript{2}; video generation remains much higher. Prefer smaller models when they suffice~\parencite{bender2021stochastic}.
      \end{block}
    \end{column}
    \hfill
    \begin{column}{0.32\textwidth}
      \begin{block}{Labor Concerns}
        Content moderation and RLHF data labeling rely on low-paid workers in the Global South. Automated code generation affects junior developer hiring pipelines.
      \end{block}
    \end{column}
    \hfill
    \begin{column}{0.32\textwidth}
      \begin{block}{Institutional Policies}
        Universities and funding agencies are developing AI use policies. Engage with your institution's process. As researchers, we have a role in shaping norms~\parencite{bommasani2021foundation}.
      \end{block}
    \end{column}
  \end{columns}
\end{frame}

% ============================================================
%  Section 10: 2026 Mini-Recap: The State of the Art
% ============================================================

\section{2026 Recap: State of the Art}

% ---- Slide 1: 2026 Update: Models --------------------------
\begin{frame}{2026 Update: Frontier Models}
  \begin{columns}[T, onlytextwidth]
    \begin{column}{0.48\textwidth}
      \textbf{Key Flagships}
      \begin{itemize}
        \item \tool{GPT-5.5}: Flagship anchor API; advanced reasoning configurations.
        \item \tool{Claude Opus 4.8}: Unmatched instruction following and extended thinking budget.
        \item \tool{Claude Sonnet 4.6}: The current gold standard for balanced speed, coding, and context.
        \item \tool{Gemini 3}: Native multimodal processing over extreme-context inputs.
      \end{itemize}
    \end{column}
    \hfill
    \begin{column}{0.48\textwidth}
      \textbf{Open-Weight Challengers}
      \begin{itemize}
        \item \tool{Gemma 4}: Designed for high-performance localized agent loops.
        \item \tool{Kimi K2}: 1T-param mixture of experts (32B active parameters).
        \item \tool{Qwen3.5}: Advanced 1M-context windows with native tool interaction.
        \item \tool{Mistral Large 3}: Heavyweight open MoE model.
      \end{itemize}
    \end{column}
  \end{columns}
  \vspace{0.6em}
  \centering\tiny\textit{Note: All models are current as of June 2026. Verify leaderboards weekly.}
\end{frame}

% ---- Slide 2: Local AI Is Now Practical --------------------
\begin{frame}[fragile]{Local AI Is Now Practical}
  \begin{columns}[T, onlytextwidth]
    \begin{column}{0.48\textwidth}
      \textbf{Local Agent Ecosystem}
      \begin{itemize}
        \item \tool{Gemma 4}: Laptop-optimized models executing native vision \& audio.
        \item \tool{Llama 4 / Kimi K2}: Heavyweight local reasoning engines.
        \item \tool{Mistral / Qwen}: Excellent local structured output.
      \end{itemize}
      \vspace{0.4em}
      \textbf{Local Tooling Pipelines}
      \begin{itemize}
        \item \tool{Ollama}: Standard local containerization.
        \item \tool{LM Studio}: User-friendly desktop server.
        \item \tool{vLLM} / \tool{llama.cpp}: Extreme performance localized inference.
      \end{itemize}
    \end{column}
    \hfill
    \begin{column}{0.48\textwidth}
      \textbf{Typical Running Setup}
      \begin{lstlisting}[language=bash, numbers=none, basicstyle=\ttfamily\tiny]
# Pull Gemma 4 edge agent
ollama pull gemma4:12b

# Start private API endpoint
vllm serve qwen3.5-plus --port 8000
      \end{lstlisting}
      \vspace{0.3em}
      \begin{block}{The Local Edge}
        Allows air-gapped drafting, private codebase refactoring, and zero running API cost on specialized academic hardware.
      \end{block}
    \end{column}
  \end{columns}
\end{frame}

% ---- Slide 3: Harness Engineering --------------------------
\begin{frame}{Harness Engineering: The New Agent Paradigm}
  \begin{columns}[T, onlytextwidth]
    \begin{column}{0.48\textwidth}
      \textbf{What is the Harness?}
      \begin{itemize}
        \item The runtime harness orchestrates the model's interaction with tools, local state, and human feedback.
        \item Shift: \textbf{Prompt} $\to$ \textbf{Context} $\to$ \textbf{Harness Engineering}.
      \end{itemize}
      \vspace{0.4em}
      \textbf{Core Architectural Components}
      $$\text{Agent} = \text{Model} + \text{Harness} + \text{Tools} + \text{Memory}$$
    \end{column}
    \hfill
    \begin{column}{0.48\textwidth}
      \textbf{Emerging Harness Options}
      \begin{itemize}
        \item \tool{Codex CLI}: Standard OpenAI task loops.
        \item \tool{LangGraph}: Graph-based checkpointing \& durable executions.
        \item \tool{Pi}: Developer-focused hackable minimalistic loops.
        \item \tool{OpenAI Agents SDK}: High-tracing hosted agents.
        \item \tool{smolagents}: Minimal Hugging Face loops.
      \end{itemize}
    \end{column}
  \end{columns}
\end{frame}

% ---- Slide 4: Persistent Agents and Memory -----------------
\begin{frame}{Persistent Agents \& Memory Architectures}
  \begin{columns}[T, onlytextwidth]
    \begin{column}{0.48\textwidth}
      \textbf{Continuous Running Agents}
      \begin{itemize}
        \item \tool{OpenClaw}: Self-hosted chat-connected inbox, files, and schedule agent.
        \item \tool{Hermes}: Memory-augmenting personal automation agent.
      \end{itemize}
      \vspace{0.4em}
      \textbf{Structured System Memory}
      \begin{itemize}
        \item \tool{Beads}: Persistent issue-tracking graphs.
        \item \textbf{Skills Repository}: Growing bank of sub-skills compiled from past runs.
      \end{itemize}
    \end{column}
    \hfill
    \begin{column}{0.48\textwidth}
      \begin{alertblock}{Persistent Memory Risks}
        \begin{itemize}
          \item \textbf{Prompt-Injection Persistence}: Poisoning memory to exploit subsequent runs.
          \item \textbf{Privacy Gaps}: Accidental reuse or upload of sensitive earlier decisions.
          \item \textbf{Stale Assumptions}: Codebase maps becoming mismatched with live files.
        \end{itemize}
      \end{alertblock}
    \end{column}
  \end{columns}
\end{frame}

% ---- Slide 5: Updated Cost & Governance --------------------
\begin{frame}{Updated Cost \& Governance}
  \begin{columns}[T, onlytextwidth]
    \begin{column}{0.48\textwidth}
      \textbf{Current Pricing Frameworks}
      \begin{itemize}
        \item Commercial flags: GPT-5.5 (\$5/\$30), Sonnet 4.6 (\$3/\$15).
        \item Local: Zero ongoing costs (only hardware amortisation).
        \item Free Tiers: Capped GitHub Copilot Free and Cursor Hobby plans.
      \end{itemize}
      \vspace{0.4em}
      \textbf{Token Optimization}
      \begin{itemize}
        \item Context compaction / summarizing old steps.
        \item Setting rigid execution budget limits.
      \end{itemize}
    \end{column}
    \hfill
    \begin{column}{0.48\textwidth}
      \textbf{Academic Governance}
      \begin{itemize}
        \item Strict data privacy policies regarding unpublished manuscripts.
        \item Human approval gates for any file write / git action.
        \item Transparent execution \& audit logs of all run loops.
      \end{itemize}
      \vspace{0.4em}
      \begin{block}{Governance Rule}
        Auditability is paramount. Log all prompt and model traces for reproducibility.
      \end{block}
    \end{column}
  \end{columns}
\end{frame}

% ============================================================
%  Section 9: Conclusion
% ============================================================

\section{Conclusion}

% ---- Key takeaways ----------------------------------------
\begin{frame}{Key Takeaways}
  \begin{enumerate}
    \item \textbf{Right tool for the task}: Inline completion, chat, and agentic tools serve different needs. Match capability to workflow rather than defaulting to one tool for everything.

    \vspace{0.4em}
    \item \textbf{Agentic oversight is mandatory}: Agents are powerful but unreliable on consequential, irreversible actions. Build in human checkpoints.

    \vspace{0.4em}
    \item \textbf{Literature tools amplify, don't replace}: Elicit, ResearchRabbit, and Semantic Scholar accelerate discovery; human judgment determines relevance and quality.

    \vspace{0.4em}
    \item \textbf{You remain the author}: AI writing assistance is legitimate; AI authorship is not. Use tools to improve your writing, not to replace your thinking.

    \vspace{0.4em}
    \item \textbf{Hallucination vigilance}: Verify every citation, every factual claim, every code snippet. Fluency is not correctness.
  \end{enumerate}
\end{frame}

% ---- Researcher's heuristic --------------------------------
\begin{frame}{A Researcher's Heuristic}
  \centering
  \begin{tikzpicture}[scale=0.95]
    % Axes
    \draw[-{Stealth[length=6pt]}, thick] (-0.3,0) -- (8.3,0)
      node[right, font=\small] {Task complexity $\rightarrow$};
    \draw[-{Stealth[length=6pt]}, thick] (0,-0.3) -- (0,5.8)
      node[above, font=\small] {Data sensitivity $\uparrow$};

    % Quadrant boxes
    \fill[blue!10]   (0.1,0.1) rectangle (4.0,2.8);
    \fill[green!10]  (4.1,0.1) rectangle (8.0,2.8);
    \fill[orange!15] (0.1,2.9) rectangle (4.0,5.6);
    \fill[red!10]    (4.1,2.9) rectangle (8.0,5.6);

    % Quadrant labels
    \node[align=center, font=\small] at (2.0, 1.6)
      {\textbf{Low sensitivity,}\\Low complexity:\\\tool{Claude} / \tool{Codex}\\Chat};
    \node[align=center, font=\small] at (6.0, 1.6)
      {\textbf{Low sensitivity,}\\High complexity:\\\tool{Claude Code} /\\Agentic workflow};
    \node[align=center, font=\small] at (2.0, 4.2)
      {\textbf{High sensitivity,}\\Low complexity:\\Local \tool{Ollama}\\(small model)};
    \node[align=center, font=\small] at (6.0, 4.2)
      {\textbf{High sensitivity,}\\High complexity:\\Local \tool{vLLM}\\(70B+ model)};

    % Midlines
    \draw[dashed, gray] (4.0, 0) -- (4.0, 5.8);
    \draw[dashed, gray] (0, 2.8) -- (8.0, 2.8);
  \end{tikzpicture}
\end{frame}

% ---- Resources --------------------------------------------
\begin{frame}{Resources \& Further Reading}
  \begin{columns}[T, onlytextwidth]
    \begin{column}{0.55\textwidth}
      \textbf{Model documentation}\\[0.2em]
      \begin{itemize}
        \item Anthropic Claude docs: \href{https://docs.anthropic.com}{docs.anthropic.com}
        \item OpenAI: \href{https://platform.openai.com/docs}{platform.openai.com/docs}
        \item Hugging Face model hub: \href{https://huggingface.co/models}{huggingface.co/models}
      \end{itemize}
      \vspace{0.5em}
      \textbf{Key papers}\\[0.2em]
      \begin{itemize}
        \item ReAct: \parencite{yao2023react}
        \item Chain-of-Thought: \parencite{wei2022cot}
        \item SWE-Bench: \parencite{jimenez2024swebench}
        \item Foundation Model risks: \parencite{bommasani2021foundation}
        \item Hallucination survey: \parencite{huang2023hallucination}
      \end{itemize}
    \end{column}
    \hfill
    \begin{column}{0.42\textwidth}
      \textbf{Tools covered today}\\[0.2em]
      \begin{itemize}
        \item MCP spec: \href{https://modelcontextprotocol.io}{modelcontextprotocol.io}
        \item Claude Code: \href{https://claude.ai/code}{claude.ai/code}
        \item Elicit: \href{https://elicit.com}{elicit.com}
        \item ResearchRabbit: \href{https://www.researchrabbit.ai}{researchrabbit.ai}
        \item Semantic Scholar API: \href{https://api.semanticscholar.org}{api.semanticscholar.org}
        \item Ollama: \href{https://ollama.ai}{ollama.ai}
      \end{itemize}
      \vspace{0.4em}
      \centering
      \qrcode[height=1.2cm]{https://github.com/your-repo/ai-tools-for-research}\\
      \tiny Slides repository
    \end{column}
  \end{columns}
\end{frame}

% ---- Call to action ----------------------------------------
\begin{frame}{Three Concrete Next Steps}
  \begin{enumerate}
    \item \textbf{Today:} Install \tool{Claude Code} or \tool{VS-Codex} and use it for your next coding task. Pay attention to where it helps and where it errs.

    \vspace{0.6em}
    \item \textbf{This week:} Set up a \tool{ResearchRabbit} collection with 3 seed papers from your current project. Explore what it surfaces.

    \vspace{0.6em}
    \item \textbf{This month:} For your next paper submission, draft one section with AI assistance. Document your process — what prompts you used, what you changed, what you verified.
  \end{enumerate}

  \vspace{0.8em}
  \begin{center}
    \small\textit{The goal is not to use AI everywhere. It is to use it \textbf{deliberately} where it genuinely accelerates your research.}
  \end{center}
\end{frame}

% ---- Questions slide ----------------------------------------
\begin{frame}[standout]
  \Large Questions \& Discussion

  \vspace{1.5em}
  \normalsize
  David Adams\\
  \href{mailto:david.adams1@student.unimelb.edu.au}{david.adams1@student.unimelb.edu.au}

  \vspace{1.2em}
  \qrcode[height=1.5cm]{https://github.com/dadams2/ai_tools_for_research}\\[0.4em]
  \small Slides available at the link above
\end{frame}


% ===========================================================
%  Appendix
% ===========================================================
\appendix

\begin{frame}{Appendix: Detailed Model Comparison (As of June 2026)}
  \small
  \begin{tabularx}{\textwidth}{lYYYY}
    \toprule
    \textbf{Model} & \textbf{Context} & \textbf{Multimodal} & \textbf{API} & \textbf{Open weight} \\
    \midrule
    \model{claude-opus-4.8} & 1M & Vision & Yes & No \\
    \model{claude-sonnet-4.6} & 1M & Vision & Yes & No \\
    \model{gpt-5.5}        & 400K & Vision+Audio & Yes & No \\
    \model{gemini-3}       & 1M+  & Vision+Audio+Video & Yes & No \\
    \model{llama-4-scout}  & 10M  & Vision & Via providers & Yes \\
    \model{mistral-large-3} & 256K & Vision & Yes & Yes \\
    \model{qwen-3.5}       & 1M   & Vision & Via providers & Yes \\
    \model{deepseek-r1}    & 128K & Text & Yes & Yes \\
    \model{gemma-4-12b}    & 128K & Vision+Audio & Via providers & Yes \\
    \bottomrule
  \end{tabularx}
\end{frame}

\begin{frame}{Appendix: Prompt Engineering Reference}
  \small
  \begin{tabularx}{\textwidth}{lX}
    \toprule
    \textbf{Pattern} & \textbf{Template} \\
    \midrule
    Role prompt & \textit{``You are an expert in [domain]. [Task]''} \\
    Reasoning / Rationale & \textit{``Use a reasoning model / request a concise rationale. [Problem]''} \\
    Few-shot & Provide 2--3 examples before the target \\
    Critique \& revise & \textit{``Review your answer for errors, then revise.''} \\
    Decompose & \textit{``Break [task] into subtasks, then solve each.''} \\
    Persona constraint & \textit{``Respond as a skeptical peer reviewer.''} \\
    Format specification & \textit{``Output a JSON object with keys: \ldots''} \\
    Context injection & Paste full paper / code before asking questions \\
    \bottomrule
  \end{tabularx}
\end{frame}

\begin{frame}{Appendix: MCP Architecture Deep-Dive}
  \centering
  \begin{tikzpicture}[
    node distance=1.2cm and 2cm,
    box/.style={draw, rounded corners=3pt, minimum width=2.6cm,
                minimum height=0.8cm, align=center, font=\small},
    arr/.style={-{Stealth[length=5pt]}, thick},
  ]
    \node[box, fill=mLightBrown!20] (host)   {Host Application\\(Claude Code)};
    \node[box, fill=blue!15, right=of host]  (client) {MCP Client};
    \node[box, fill=green!15, right=of client](server) {MCP Server};
    \node[box, fill=gray!15, above right=0.6cm and 0.3cm of server] (res1) {Resources};
    \node[box, fill=gray!15, right=0.3cm of server]                 (res2) {Tools};
    \node[box, fill=gray!15, below right=0.6cm and 0.3cm of server] (res3) {Prompts};

    \draw[arr] (host.east)   -- (client.west)  node[midway, above, font=\tiny]{JSON-RPC 2.0};
    \draw[arr] (client.east) -- (server.west)  node[midway, above, font=\tiny]{stdio/SSE};
    \draw[arr] (server.north east) -- (res1.west);
    \draw[arr] (server.east)       -- (res2.west);
    \draw[arr] (server.south east) -- (res3.west);
    \draw[arr, dashed] (res1.west) -- (server.north east);
    \draw[arr, dashed] (res2.west) -- (server.east);
    \draw[arr, dashed] (res3.west) -- (server.south east);
  \end{tikzpicture}

  \vspace{0.5em}
  \small Servers expose \emph{resources} (data), \emph{tools} (actions), and \emph{prompts} (templates).\\
  Clients \& servers communicate over JSON-RPC 2.0 via stdio or SSE transport.
\end{frame}

\begin{frame}[fragile]{Appendix: Claude Code via GitHub Copilot Licence}
  \scriptsize
  Use your \textbf{academic GitHub Copilot licence} as a free model provider for Claude Code.\\[0.2em]
  \url{https://github.com/feiskyer/claude-code-settings/blob/main/guidances/github-copilot.md}

  \smallskip
  \textbf{Step 1 — Install}
  \begin{lstlisting}[language=bash, numbers=none, basicstyle=\ttfamily\tiny, aboveskip=1pt, belowskip=2pt]
npm install -g copilot-api @anthropic-ai/claude-code
  \end{lstlisting}

  \textbf{Step 2 — Start the proxy \& authenticate}
  \begin{lstlisting}[language=bash, numbers=none, basicstyle=\ttfamily\tiny, aboveskip=1pt, belowskip=2pt]
copilot-api start --proxy-env
# Visit https://github.com/login/device and enter the shown code
  \end{lstlisting}

  \textbf{Step 3 — Create \texttt{\textasciitilde/.claude/settings.json}}
  \begin{lstlisting}[language=json, numbers=none, basicstyle=\ttfamily\tiny, aboveskip=1pt, belowskip=2pt]
{
  "env": {
    "ANTHROPIC_BASE_URL": "http://localhost:4141",
    "ANTHROPIC_AUTH_TOKEN": "sk-dummy",
    "ANTHROPIC_MODEL": "claude-sonnet-4.5",
    "DISABLE_NON_ESSENTIAL_MODEL_CALLS": "1"
  }
}
  \end{lstlisting}

  \textbf{Step 4 — Run Claude Code}\\[0.1em]
  Open a new terminal and run \texttt{claude}.
  Available models include \model{claude-opus-4}, \model{gemini-2.5-pro}, \model{o3}, and more.
\end{frame}

\begin{frame}[allowframebreaks]{References}
  \printbibliography[heading=none]
\end{frame}

\end{document}
